\chapter{Stata / R 命令速查}
\label{ch:UappB-stata}

本附录把正文各章用到的 Stata 与 R 命令按主题汇总为“命令 $|$ 作用 $|$ 常用选项 / 注意 $|$ 章”的速查表，并在最后一节给出三段完整的实证流程模板和一份常见陷阱清单。考试不要求上机，但题目可能问“如何估计”“用什么命令或选项”，作答时写出核心命令与关键选项即可。表中“章”一栏是该方法在正文中所在的章。

\begin{summary}[标记约定与 Stata 语法要点]
\begin{itemize}
  \item {[SSC]}：用户编写的命令，首次使用前运行 \texttt{ssc install 命令名}；{[官方]}：Stata 官方命令。
  \item \tagext：正文未展开、为方便查找而列出的常用写法（均为标准语法）。
  \item 因子变量：\texttt{i.x} 为分类变量，\texttt{c.x} 为连续变量；\texttt{\#} 只生成交互项，\texttt{\#\#} 生成主效应 + 交互项。
  \item 时间序列算子：\texttt{L.}、\texttt{F.}、\texttt{D.} 分别为滞后、超前、一阶差分（\texttt{LD.y} 即 $\Delta y_{t-1}$）；使用前必须先 \texttt{xtset} 或 \texttt{tsset}。
  \item 方差估计：\texttt{vce(robust)} 为异方差稳健，\texttt{vce(cluster id)} 为按 \texttt{id} 聚类；多数命令也接受 \texttt{robust}（简写 \texttt{r}）作为 \texttt{vce(robust)} 的同义写法。
\end{itemize}
\end{summary}

%======================================================================
\section{OLS、稳健/聚类标准误与假设检验（第 4–6 章）}\label{sec:UappB-ols}

\subsection{估计与标准误：同方差、HC、聚类、双向聚类与 HAC}

{\renewcommand{\arraystretch}{1.12}\begin{longtable}{@{}>{\raggedright\arraybackslash\small}p{5.3cm}>{\raggedright\arraybackslash\small}p{3.8cm}>{\raggedright\arraybackslash\small}p{5.3cm}>{\centering\arraybackslash\small}p{0.9cm}@{}}
\toprule
\textbf{命令} & \textbf{作用} & \textbf{常用选项 / 注意} & \textbf{章}\\
\midrule
\endfirsthead
\toprule
\textbf{命令} & \textbf{作用} & \textbf{常用选项 / 注意} & \textbf{章}\\
\midrule
\endhead
\bottomrule
\endlastfoot

\texttt{regress y x1 x2} & OLS（简写 \texttt{reg}）；默认同方差标准误 $s^2(X\T X)\inv$ & 自动含截距；\texttt{noconstant} 去截距 & 4–5\\
\addlinespace[1pt]
\texttt{reg y x1 x2, vce(robust)} & HC1 异方差稳健标准误 & 等价写法 \texttt{robust}、\texttt{r}；\texttt{vce(hc2)}、\texttt{vce(hc3)} 为 HC2、HC3 & 5\\
\addlinespace[1pt]
\texttt{reg y x1 x2, vce(cluster firmid)} & 聚类稳健标准误（CRVE） & 等价 \texttt{cluster(firmid)}；方差矩阵乘小样本调整因子 $\frac{n-1}{n-K}\cdot\frac{G}{G-1}$，临界值用 $t_{G-1}$（$G$ 为簇数）；点估计不变 & 5\\
\addlinespace[1pt]
\texttt{reghdfe y x, absorb(firm year) vce(cluster firm year)}\newline {\footnotesize\rmfamily [SSC]} & 双向聚类标准误 $\hat V_A+\hat V_B-\hat V_{A\cap B}$ & 两个维度的簇数都要足够多 & 5\\
\addlinespace[1pt]
\texttt{tsset t}\newline \texttt{newey y x, lag(4)} & HAC（Newey–West）标准误 & \texttt{lag()} 为截断滞后 $L$，常取 $0.75\,T^{1/3}$ & 5\\
\addlinespace[1pt]
\texttt{estat vif} & 方差膨胀因子 $\mathrm{VIF}_j$ & 在 \texttt{regress} 之后使用；$\mathrm{VIF}>10$ 提示严重共线 & 3\\
\addlinespace[1pt]
\texttt{predict uhat, residuals}\newline \texttt{predict yhat, xb}\newline \tagext & 生成残差 / 拟合值 & 用于诊断、自建检验 & —\\
\addlinespace[1pt]
\texttt{estimates store m1}\newline \texttt{estimates table m1 m2, b se} & 保存结果并列比较 & \texttt{keep(x \_cons)} 只显示指定系数 & 4, 10\\
\end{longtable}}


\subsection{假设检验：Wald/\texorpdfstring{$F$}{F}、线性组合与 Delta 方法}

{\renewcommand{\arraystretch}{1.12}\begin{longtable}{@{}>{\raggedright\arraybackslash\small}p{5.3cm}>{\raggedright\arraybackslash\small}p{3.8cm}>{\raggedright\arraybackslash\small}p{5.3cm}>{\centering\arraybackslash\small}p{0.9cm}@{}}
\toprule
\textbf{命令} & \textbf{作用} & \textbf{常用选项 / 注意} & \textbf{章}\\
\midrule
\endfirsthead
\toprule
\textbf{命令} & \textbf{作用} & \textbf{常用选项 / 注意} & \textbf{章}\\
\midrule
\endhead
\bottomrule
\endlastfoot

\texttt{test x1 x2} & Wald/$F$ 检验 $H_0:\beta_1=\beta_2=0$（$q=2$） & 使用回归时指定的方差（稳健/聚类）；报告 $F=W/q$ & 6\\
\addlinespace[1pt]
\texttt{test (x1 = 1) (x2 = 1)} & $H_0:\beta_1=1$ 且 $\beta_2=1$ & 每个约束放一对括号，一条命令内为联合检验 & 6\\
\addlinespace[1pt]
\texttt{test x1 + 2*x2 = 0}\newline \texttt{test 2*x1 + 3*x2 = 4} & 一般线性约束 $R\beta=r$：分别检验 $H_0:\beta_1+2\beta_2=0$ 与 $H_0:2\beta_1+3\beta_2=4$ & 两行是两个\key{单独}的检验（各 $q=1$）。要联合检验（$q=2$），写 \texttt{test (x1 + 2*x2 = 0) (2*x1 + 3*x2 = 4)}，或在第二行末加 \texttt{, accumulate} & 6\\
\addlinespace[1pt]
\texttt{testparm x1-x3}\newline \texttt{testparm i.year}\newline \tagext & 一组变量（含因子变量）的系数全为 0 & 常用于年份虚拟变量、事前系数的联合检验 & 6\\
\addlinespace[1pt]
\texttt{lincom x1 + 2*x2} & 单个线性组合的估计值、se、$t$ 与置信区间 & 只针对一个线性组合 & 6\\
\addlinespace[1pt]
\texttt{testnl \_b[x1]/\_b[x2] = 1} & 非线性约束的 Wald 检验 & 结果依赖约束的写法：$\beta_1/\beta_2=1$ 与 $\beta_1-\beta_2=0$ 的 Wald 值不同 & 6\\
\addlinespace[1pt]
\texttt{nlcom -\_b[exper]/(2*\_b[exper2])} & 非线性函数的 Delta 方法 se & 如工资–工龄曲线的拐点 $-\beta_1/(2\beta_2)$；若平方项用因子变量写成 \texttt{c.exper\#c.exper}，系数名为 \texttt{\_b[c.exper\#c.exper]} & 5–6\\
\end{longtable}}


\begin{english}
\begin{itemize}
  \item We estimate the model by OLS and report heteroskedasticity-robust (HC1) standard errors.
  \item Standard errors are clustered at the firm level, which allows arbitrary correlation of the errors within a firm; with $G$ clusters, inference is based on the $t_{G-1}$ distribution.
  \item The joint null hypothesis $H_0:\beta_1=\beta_2=0$ is tested by a Wald test based on the robust covariance matrix and reported as $F=W/q$.
  \item The standard error of the turning point $-\beta_1/(2\beta_2)$ is obtained by the delta method.
\end{itemize}
\end{english}

%======================================================================
\section{IV、2SLS、GMM 与三大检验（第 8–9 章）}\label{sec:UappB-iv}

\subsection{估计命令：2SLS、两步 GMM、LIML 与 ivreg2}

{\renewcommand{\arraystretch}{1.12}\begin{longtable}{@{}>{\raggedright\arraybackslash\small}p{5.3cm}>{\raggedright\arraybackslash\small}p{3.8cm}>{\raggedright\arraybackslash\small}p{5.3cm}>{\centering\arraybackslash\small}p{0.9cm}@{}}
\toprule
\textbf{命令} & \textbf{作用} & \textbf{常用选项 / 注意} & \textbf{章}\\
\midrule
\endfirsthead
\toprule
\textbf{命令} & \textbf{作用} & \textbf{常用选项 / 注意} & \textbf{章}\\
\midrule
\endhead
\bottomrule
\endlastfoot

\texttt{ivregress 2sls y (x1 = z1 z2) x2, vce(robust) first}\newline {\footnotesize\rmfamily [官方]} & 2SLS：$x_1$ 内生，$z_1,z_2$ 为排除工具，$x_2$ 外生（自己作自己的工具） & \texttt{vce(robust)} 稳健 se；\texttt{first} 打印第一阶段；\texttt{vce(cluster id)} 聚类 se & 8–9\\
\addlinespace[1pt]
\texttt{ivregress gmm y (x1 = z1 z2) x2, wmatrix(robust)} & 两步有效 GMM（异方差稳健权重矩阵） & \texttt{igmm} 为迭代 GMM；之后 \texttt{estat overid} 报告 Hansen $J$ & 8\\
\addlinespace[1pt]
\texttt{ivregress liml y (x1 = z1 z2) x2, vce(robust)} & LIML，弱工具时的替代 & 工具多时比 2SLS 稳健 & 8–9\\
\addlinespace[1pt]
\texttt{ivreg2 y x2 (x1 = z1 z2), robust first}\newline {\footnotesize\rmfamily [SSC]} & 报告 Kleibergen–Paap rk Wald $F$ 等弱工具诊断 & 异方差/聚类稳健的弱工具检验；\texttt{gmm2s} 为两步 GMM，\texttt{liml} 为 LIML，\texttt{cluster(id)} 聚类，\texttt{endog(x1)} 检验 $x_1$ 的内生性 & 9\\
\end{longtable}}


\subsection{估计后检验：相关性、过度识别与内生性}

{\renewcommand{\arraystretch}{1.12}\begin{longtable}{@{}>{\raggedright\arraybackslash\small}p{5.3cm}>{\raggedright\arraybackslash\small}p{3.8cm}>{\raggedright\arraybackslash\small}p{5.3cm}>{\centering\arraybackslash\small}p{0.9cm}@{}}
\toprule
\textbf{命令} & \textbf{作用} & \textbf{常用选项 / 注意} & \textbf{章}\\
\midrule
\endfirsthead
\toprule
\textbf{命令} & \textbf{作用} & \textbf{常用选项 / 注意} & \textbf{章}\\
\midrule
\endhead
\bottomrule
\endlastfoot

\texttt{estat firststage} & 检验一（相关性）：第一阶段 $F$、偏 $R^2$ & 非稳健设定下还报告 Cragg–Donald 最小特征值统计量与 Stock–Yogo 临界值 & 9\\
\addlinespace[1pt]
\texttt{estat overid} & 检验二（过度识别）：Sargan、Basmann；稳健设定下为稳健 score 检验；GMM 之后为 Hansen $J$ & 只在 $L>K$ 时可做 & 9\\
\addlinespace[1pt]
\texttt{estat endogenous} & 检验三（内生性） & 可简写 \texttt{estat endog}。非稳健：Durbin $\chi^2$、Wu–Hausman $F$；稳健：稳健 score $\chi^2$、稳健回归 $F$ & 9\\
\addlinespace[1pt]
\texttt{regress y x1 x2}\newline \texttt{estimates store ols}\newline \texttt{ivregress 2sls y (x1 = z1 z2) x2}\newline \texttt{estimates store iv}\newline \texttt{hausman iv ols} & 传统 Hausman 检验（同方差） & 顺序：先一致的（iv），后有效的（ols）；自由度应为内生变量个数，可用 \texttt{df(1)} 指定 & 9\\
\addlinespace[1pt]
\texttt{test}、\texttt{testnl}（\texttt{ivregress} 之后） & 2SLS/GMM 系数的 Wald 检验 & 自动使用估计时指定的方差估计 & 8\\
\end{longtable}}


\subsection{2SLS 不要手工两步做}

\begin{note}[手工两步 2SLS 的标准误是错的]
手工“先把 $x_1$ 对全部工具回归取拟合值 $\hat x_1$，再把 $y$ 对 $\hat x_1$ 与 $x_2$ 回归”得到的\key{系数正确}，但第二步软件计算标准误时用的残差是 $y_i-\hat x_i\T\hat\beta$（$\hat x_i$ 为拟合值），\key{是错的}。正确的残差必须用原始回归元：$\hat u_i=y_i-x_i\T\hat\beta_{2SLS}$。所以应直接用 \texttt{ivregress}（第 8 章）。
\end{note}

\begin{english}
\begin{itemize}
  \item We estimate the equation by 2SLS, treating $x_1$ as endogenous and using $z_1$ and $z_2$ as excluded instruments; the exogenous regressor $x_2$ serves as its own instrument.
  \item Instrument relevance is assessed by the first-stage $F$ statistic, the overidentifying restrictions by the Sargan/Hansen $J$ test, and the endogeneity of $x_1$ by the Durbin--Wu--Hausman test.
  \item 2SLS should not be computed by two manual OLS steps, because the second-step standard errors are based on residuals computed with the fitted values instead of the original regressors.
\end{itemize}
\end{english}

%======================================================================
\section{静态面板（第 10–11 章）}\label{sec:UappB-static}

\subsection{固定效应、LSDV 与双向固定效应}

{\renewcommand{\arraystretch}{1.12}\begin{longtable}{@{}>{\raggedright\arraybackslash\small}p{5.3cm}>{\raggedright\arraybackslash\small}p{3.8cm}>{\raggedright\arraybackslash\small}p{5.3cm}>{\centering\arraybackslash\small}p{0.9cm}@{}}
\toprule
\textbf{命令} & \textbf{作用} & \textbf{常用选项 / 注意} & \textbf{章}\\
\midrule
\endfirsthead
\toprule
\textbf{命令} & \textbf{作用} & \textbf{常用选项 / 注意} & \textbf{章}\\
\midrule
\endhead
\bottomrule
\endlastfoot

\texttt{xtset id year}\newline {\footnotesize\rmfamily [官方]} & 声明面板结构 & 所有 \texttt{xt} 命令与 \texttt{L.}、\texttt{D.} 算子之前必须先做 & 10\\
\addlinespace[1pt]
\texttt{reg y x i.id, vce(cluster id)} & LSDV：OLS + 全部个体虚拟变量 & 系数同 FE；\texttt{estimates table, keep(x) b se} 只看斜率 & 10\\
\addlinespace[1pt]
\texttt{xtreg y x, fe} & WG/FE 估计，经典 se（自由度 $NT-N-K$） & 报告 \texttt{sigma\_u}、\texttt{sigma\_e}、\texttt{rho} 与“所有 $\alpha_i=0$”的 $F$ 检验；Stata 的 \texttt{sigma\_u} 是本书的 $\sigma_\alpha$，\texttt{sigma\_e} 是本书的 $\sigma_u$ & 10\\
\addlinespace[1pt]
\texttt{xtreg y x, fe vce(cluster id)} & FE + 按个体聚类 se（Arellano） & 对 \texttt{xtreg} 而言 \texttt{vce(robust)} 也是按个体聚类 & 10\\
\addlinespace[1pt]
\texttt{areg y x, absorb(id) vce(cluster id)} & 吸收个体效应 & 系数同 FE；聚类 se 的自由度调整把被吸收的个体虚拟变量计入，短面板中 se 偏大 & 10\\
\addlinespace[1pt]
\texttt{xtreg y x i.year, fe vce(cluster id)} & TWFE：个体 FE + 年份虚拟变量 & 平衡与非平衡面板都正确（虚拟变量法） & 10\\
\addlinespace[1pt]
\texttt{reg y x i.id i.year, vce(cluster id)} & 双向 LSDV & 与上式系数相同 & 10\\
\addlinespace[1pt]
\texttt{reghdfe y x, absorb(id year) vce(cluster id)}\newline {\footnotesize\rmfamily [SSC]} & 高维多向固定效应 & \texttt{absorb(firm city\#year)} 吸收交互效应；速度快 & 10\\
\addlinespace[1pt]
\texttt{xtset city year}\newline \texttt{xtreg y x i.province, fe vce(robust)} & 嵌套多层效应：省份效应被城市效应吸收 & 每个城市只属于一个省，\texttt{i.province} 与城市固定效应完全共线而被自动略去；FE 无法识别更高层级的不随时间变化的效应 & 10\\
\end{longtable}}


\subsection{随机效应、Hausman 检验、FEF 与 Mundlak 回归}

{\renewcommand{\arraystretch}{1.12}\begin{longtable}{@{}>{\raggedright\arraybackslash\small}p{5.3cm}>{\raggedright\arraybackslash\small}p{3.8cm}>{\raggedright\arraybackslash\small}p{5.3cm}>{\centering\arraybackslash\small}p{0.9cm}@{}}
\toprule
\textbf{命令} & \textbf{作用} & \textbf{常用选项 / 注意} & \textbf{章}\\
\midrule
\endfirsthead
\toprule
\textbf{命令} & \textbf{作用} & \textbf{常用选项 / 注意} & \textbf{章}\\
\midrule
\endhead
\bottomrule
\endlastfoot

\texttt{xtreg y x z, re} & RE-FGLS（准去均值）；可含时不变变量 $z_i$ & \texttt{theta} 报告 $\hat\theta$；\texttt{vce(cluster id)} 聚类 se & 11\\
\addlinespace[1pt]
\texttt{xtreg y x, be} & 组间估计量（个体均值回归） & RE 是 WG 与组间估计量的加权平均 & 11\\
\addlinespace[1pt]
\texttt{xtreg y x1 x2, fe}\newline \texttt{estimates store fe}\newline \texttt{xtreg y x1 x2, re}\newline \texttt{estimates store re}\newline \texttt{hausman fe re, sigmamore} & Hausman 检验 FE vs RE（$\chi^2_K$） & 先一致的（fe）后有效的（re）；\texttt{sigmamore} 让两个方差都用 RE 的 $\hat\sigma^2$，避免方差之差非正定 & 11\\
\addlinespace[1pt]
\texttt{xtreg y x, fe}\newline \texttt{gen r = y - \_b[x]*x}\newline \texttt{bysort id: egen rbar = mean(r)}\newline \texttt{bysort id: gen first = (\_n==1)}\newline \texttt{reg rbar z if first} & FEF 两步法：点估计 $\hat\gamma$ & 第二步 se 忽略了第一步的估计误差，不能直接用 & 11\\
\addlinespace[1pt]
\texttt{bysort id: egen xbar = mean(x)}\newline \texttt{gen xdot = x - xbar}\newline \texttt{ivregress 2sls y z (x = xdot), vce(cluster id)} & FEF 的 IV 等价形式（Hansen §17.32） & 点估计与两步法相同，且给出正确的聚类 se & 11\\
\addlinespace[1pt]
\texttt{xtreg y x xbar z, re vce(cluster id)}\newline \texttt{test xbar} & Mundlak / CRE 回归；检验 $\pi=0$ & 即稳健的 Hausman 检验 & 11\\
\addlinespace[1pt]
\texttt{xtlogit y x, fe}\newline \texttt{clogit y x, group(id)} & 面板 logit 的 Chamberlain 条件似然（两条命令等价） & 只用 $y_{it}$ 在期间内发生过变化的个体；无法得到边际效应 & 11\\
\addlinespace[1pt]
\texttt{xtprobit y x, re} & RE probit（对 $\alpha_i$ 积分） & 需 $\alpha_i$ 与 $x_{it}$ 独立 & 11\\
\end{longtable}}


\begin{english}
\begin{itemize}
  \item We estimate the model by fixed effects (the within estimator) and cluster the standard errors by individual.
  \item Year dummies are included to absorb common time effects, so the specification is a two-way fixed effects model.
  \item The Hausman test compares FE and RE: under the null that $\alpha_i$ is uncorrelated with $x_{it}$ both estimators are consistent and RE is efficient, while under the alternative only FE is consistent.
  \item A cluster-robust version of the test is the Mundlak (correlated random effects) regression, which tests whether the coefficients on $\bar x_{i\cdot}$ are jointly zero.
\end{itemize}
\end{english}

%======================================================================
\section{动态面板（第 12 章）}\label{sec:UappB-dynamic}

\subsection{\texttt{xtabond2} 的三种估计量与关键选项}

\subsubsection{差分 GMM、系统 GMM 与 FOD GMM 的切换}

{\renewcommand{\arraystretch}{1.12}\begin{longtable}{@{}>{\raggedright\arraybackslash\small}p{5.3cm}>{\raggedright\arraybackslash\small}p{3.8cm}>{\raggedright\arraybackslash\small}p{5.3cm}>{\centering\arraybackslash\small}p{0.9cm}@{}}
\toprule
\textbf{命令} & \textbf{作用} & \textbf{常用选项 / 注意} & \textbf{章}\\
\midrule
\endfirsthead
\toprule
\textbf{命令} & \textbf{作用} & \textbf{常用选项 / 注意} & \textbf{章}\\
\midrule
\endhead
\bottomrule
\endlastfoot

\texttt{ssc install xtabond2}\newline \texttt{xtset id year} & 安装（Roodman 2009）并声明面板 & — & 12\\
\addlinespace[1pt]
\texttt{xtabond2 y L.y x, gmm(L.y) iv(x) noleveleq}\newline {\footnotesize\rmfamily [SSC]} & Arellano–Bond 差分 GMM & \texttt{noleveleq}：不用水平方程 & 12\\
\addlinespace[1pt]
\texttt{xtabond2 y L.y x, gmm(L.y) iv(x)} & Blundell–Bond 系统 GMM & 默认加入水平方程，额外矩条件 $\E[(y_{it}-\rho y_{i,t-1})\Delta y_{i,t-1}]=0$ & 12\\
\addlinespace[1pt]
\texttt{xtabond2 y L.y x, gmm(L.y) iv(x) noleveleq orthogonal} & FOD GMM & \texttt{orthogonal}：用前向正交离差代替一阶差分 & 12\\
\end{longtable}}


\subsubsection{工具设定与两步估计选项}

{\renewcommand{\arraystretch}{1.12}\begin{longtable}{@{}>{\raggedright\arraybackslash\small}p{5.3cm}>{\raggedright\arraybackslash\small}p{3.8cm}>{\raggedright\arraybackslash\small}p{5.3cm}>{\centering\arraybackslash\small}p{0.9cm}@{}}
\toprule
\textbf{命令} & \textbf{作用} & \textbf{常用选项 / 注意} & \textbf{章}\\
\midrule
\endfirsthead
\toprule
\textbf{命令} & \textbf{作用} & \textbf{常用选项 / 注意} & \textbf{章}\\
\midrule
\endhead
\bottomrule
\endlastfoot

\texttt{gmmstyle()}（简写 \texttt{gmm()}） & 内生/前定变量的滞后按“每个滞后、每期一列”作工具 & 即 AB 的块对角工具矩阵 $W_i$ & 12\\
\addlinespace[1pt]
\texttt{ivstyle()}（简写 \texttt{iv()}） & 外生变量按普通 IV 方式作工具（每个变量一列） & 如 \texttt{iv(x yr*)} & 12\\
\addlinespace[1pt]
\texttt{twostep robust} & 两步 GMM + Windmeijer 校正 se & \texttt{small}：报告小样本 $t$/$F$ & 12\\
\addlinespace[1pt]
\texttt{gmm(L.y, laglimits(1 3) collapse)}\newline \tagext & 限制工具数 & \texttt{laglimits(a b)} 只用第 $a$–$b$ 阶滞后；\texttt{collapse} 把每个滞后合并为一列 & 12\\
\addlinespace[1pt]
\texttt{gmm(x, laglimits(1 .))}\newline \texttt{gmm(x, laglimits(2 .))} & $x$ 为前定 / 内生时的工具 & 前定从 $x_{i,t-1}$ 起，内生从 $x_{i,t-2}$ 起；此时 $x$ 要从 \texttt{iv()} 移到 \texttt{gmm()} & 12\\
\end{longtable}}


\subsubsection{输出底部的设定检验}

{\renewcommand{\arraystretch}{1.12}\begin{longtable}{@{}>{\raggedright\arraybackslash\small}p{5.3cm}>{\raggedright\arraybackslash\small}p{3.8cm}>{\raggedright\arraybackslash\small}p{5.3cm}>{\centering\arraybackslash\small}p{0.9cm}@{}}
\toprule
\textbf{命令} & \textbf{作用} & \textbf{常用选项 / 注意} & \textbf{章}\\
\midrule
\endfirsthead
\toprule
\textbf{命令} & \textbf{作用} & \textbf{常用选项 / 注意} & \textbf{章}\\
\midrule
\endhead
\bottomrule
\endlastfoot

（输出底部，无需另写命令） & AR(1)/AR(2) 检验、Sargan、Hansen、Difference-in-Hansen & AR(1) 应拒绝、AR(2) 不应拒绝；工具数应少于个体数 & 12\\
\end{longtable}}


\subsection{官方命令、手工估计、估计量的“界”与单位根检验}

{\renewcommand{\arraystretch}{1.12}\begin{longtable}{@{}>{\raggedright\arraybackslash\small}p{5.3cm}>{\raggedright\arraybackslash\small}p{3.8cm}>{\raggedright\arraybackslash\small}p{5.3cm}>{\centering\arraybackslash\small}p{0.9cm}@{}}
\toprule
\textbf{命令} & \textbf{作用} & \textbf{常用选项 / 注意} & \textbf{章}\\
\midrule
\endfirsthead
\toprule
\textbf{命令} & \textbf{作用} & \textbf{常用选项 / 注意} & \textbf{章}\\
\midrule
\endhead
\bottomrule
\endlastfoot

\texttt{xtabond y x, lags(1) twostep vce(robust)}\newline {\footnotesize\rmfamily [官方]} & AB 差分 GMM & 之后 \texttt{estat abond} 报告 AR 检验 & 12\\
\addlinespace[1pt]
\texttt{xtdpdsys y x, lags(1) twostep vce(robust)}\newline {\footnotesize\rmfamily [官方]} & BB 系统 GMM & \texttt{xtdpd} 为通用命令 & 12\\
\addlinespace[1pt]
\texttt{reg D.y LD.y, noconstant} & FD-OLS（不一致：$\plim=(\rho-1)/2$） & 对照用 & 12\\
\addlinespace[1pt]
\texttt{ivregress 2sls D.y (LD.y = L2.y), noconstant} & Anderson–Hsiao IV & 用 $y_{i,t-2}$ 作 $\Delta y_{i,t-1}$ 的工具 & 12\\
\addlinespace[1pt]
\texttt{reg y L.y x yr*, vce(cluster id)}\newline \texttt{xtreg y L.y x yr*, fe vce(cluster id)} & 混合 OLS（上偏）、WG（下偏）两个“界” & 一致估计应落在二者之间 & 12\\
\addlinespace[1pt]
\texttt{tab year, gen(yr)} & 生成年份虚拟变量 \texttt{yr1}、\texttt{yr2}、… & 放入 \texttt{iv(yr*)} & 12\\
\addlinespace[1pt]
\texttt{xtunitroot llc y}\newline \texttt{xtunitroot ips y} & 面板单位根检验（LLC、IPS） & $T$ 较大时先做 & 12\\
\addlinespace[1pt]
\texttt{xtcointtest pedroni y x}\newline \tagext & 面板协整检验 & 另有 \texttt{xtcointtest kao}、\texttt{xtcointtest westerlund} & 12\\
\end{longtable}}


\begin{english}
\begin{itemize}
  \item We estimate the dynamic panel model by system GMM, using lags of $y_{it}$ dated $t-2$ and earlier as instruments for the differenced equation and lagged differences as instruments for the level equation.
  \item We report two-step estimates with Windmeijer-corrected standard errors.
  \item The AR(1) test rejects while the AR(2) test does not, which supports the validity of the moment conditions; the Hansen $J$ test does not reject the overidentifying restrictions.
  \item To avoid instrument proliferation we collapse the instrument matrix and report the number of instruments, which is smaller than the number of individuals.
\end{itemize}
\end{english}

%======================================================================
\section{截面相关、因子模型与交互固定效应（第 13–14 章）}\label{sec:UappB-factor}

\subsection{Stata 命令：PCA、交互固定效应、CCE 与截面相关检验}

{\renewcommand{\arraystretch}{1.12}\begin{longtable}{@{}>{\raggedright\arraybackslash\small}p{5.3cm}>{\raggedright\arraybackslash\small}p{3.8cm}>{\raggedright\arraybackslash\small}p{5.3cm}>{\centering\arraybackslash\small}p{0.9cm}@{}}
\toprule
\textbf{命令} & \textbf{作用} & \textbf{常用选项 / 注意} & \textbf{章}\\
\midrule
\endfirsthead
\toprule
\textbf{命令} & \textbf{作用} & \textbf{常用选项 / 注意} & \textbf{章}\\
\midrule
\endhead
\bottomrule
\endlastfoot

\texttt{pca y1-y20} & 统计 PCA（默认对相关矩阵） & 变量在列上，与本书 $U$（$N\times T$）互为转置 & 13\\
\addlinespace[1pt]
\texttt{factor y1-y20, pcf factors(2)}\newline \texttt{predict f1 f2} & 主成分因子法与因子得分 & Hansen §11.13–11.14；\texttt{factor y1-y20, ml factors(2)} 为极大似然 & 13\\
\addlinespace[1pt]
\texttt{regife y x1 x2, factors(id year) dimension(2)}\newline {\footnotesize\rmfamily [SSC]} & Bai (2009) 交互固定效应 & \texttt{absorb(id year)} 同时吸收双向 FE；完整语法见 \texttt{help regife} & 14\\
\addlinespace[1pt]
\texttt{xtmg y x1 x2, cce}\newline {\footnotesize\rmfamily [SSC]} & CCEMG（Eberhardt） & 不加 \texttt{cce} 为 Pesaran–Smith MG；完整语法见 \texttt{help xtmg} & 14\\
\addlinespace[1pt]
\texttt{xtdcce2 y x1 x2, cr(y x1 x2)}\newline {\footnotesize\rmfamily [SSC]} & CCE 均值组（Ditzen） & \texttt{cr()} 指定取截面平均的变量；\texttt{pooled()} 设混合系数 & 14\\
\addlinespace[1pt]
\texttt{xtreg y x1 x2, fe}\newline \texttt{xtcsd, pesaran abs}\newline {\footnotesize\rmfamily [SSC]} & FE 残差的 Pesaran CD 检验（\texttt{xtcsd} 为用户命令） & \texttt{abs} 报告平均绝对相关；\texttt{friedman}、\texttt{frees} 为其他检验 & 14\\
\addlinespace[1pt]
\texttt{predict ehat, e}\newline \texttt{xtcd ehat}\newline {\footnotesize\rmfamily [SSC]} & 对残差（或变量）做 CD 检验（Eberhardt） & \texttt{predict ..., e} 取 FE 回归的特异残差 $\hat u_{it}$ & 14\\
\addlinespace[1pt]
\texttt{xttest2}\newline {\footnotesize\rmfamily [SSC]} & Breusch–Pagan LM 截面独立检验 & 在 \texttt{xtreg, fe} 之后使用；要求 $T$ 相对 $N$ 较大 & 14\\
\addlinespace[1pt]
\texttt{xtscc y x1 x2, fe lag(2)}\newline {\footnotesize\rmfamily [SSC]} & Driscoll–Kraay 标准误（截面相关 + 序列相关） & 依赖 $T$ 大；只修正 se，不修正内生性 & 13\\
\end{longtable}}


\subsection{R 与矩阵语言：PCA、交互固定效应、CCE 与 CD 检验}

{\renewcommand{\arraystretch}{1.12}\begin{longtable}{@{}>{\raggedright\arraybackslash\small}p{5.3cm}>{\raggedright\arraybackslash\small}p{3.8cm}>{\raggedright\arraybackslash\small}p{5.3cm}>{\centering\arraybackslash\small}p{0.9cm}@{}}
\toprule
\textbf{命令} & \textbf{作用} & \textbf{常用选项 / 注意} & \textbf{章}\\
\midrule
\endfirsthead
\toprule
\textbf{命令} & \textbf{作用} & \textbf{常用选项 / 注意} & \textbf{章}\\
\midrule
\endhead
\bottomrule
\endlastfoot

\texttt{e <- eigen(t(U) \%*\% U, symmetric = TRUE)}\newline \texttt{Fhat <- e\$vectors[, 1:r]}\newline \texttt{Lhat <- U \%*\% Fhat} & PCA：$\hat F$ 为 $U\T U$ 前 $r$ 个特征向量，$\hat\Lambda=U\hat F$（标准化 $F\T F=I_r$） & R 的 \texttt{eigen} 按特征值降序返回；若标准化 $F\T F/T=I_r$，改用 \texttt{Fhat <- sqrt(T) * e\$vectors[, 1:r]} 与 \texttt{Lhat <- U \%*\% Fhat / T}；\texttt{prcomp(X, scale. = TRUE)} 为统计 PCA & 13\\
\addlinespace[1pt]
\texttt{[V, D] = eig(U\textquotesingle{}*U)}\newline \texttt{np.linalg.eigh(U.T @ U)} & 同上（第一行 MATLAB，第二行 Python） & 二者均按特征值\key{升序}返回，需重排 & 13\\
\addlinespace[1pt]
\texttt{library(phtt)}\newline \texttt{Eup(Y \textasciitilde{} X1 + X2, additive.effects = "individual", factor.dim = 2)} & Bai 型交互固定效应迭代估计 & 另有 \texttt{KSS()}；参数名以 \texttt{?Eup} 为准 & 14\\
\addlinespace[1pt]
\texttt{pd <- pdata.frame(df, index = c("id", "year"))} & 声明面板（plm 包） & — & 10, 14\\
\addlinespace[1pt]
\texttt{pcce(y \textasciitilde{} x1 + x2, data = pd, model = "p")} & PCCE（CCEP） & \texttt{model = "mg"} 为 CCEMG & 14\\
\addlinespace[1pt]
\texttt{pmg(y \textasciitilde{} x1 + x2, data = pd, model = "mg")} & Pesaran–Smith 均值组 & \texttt{"cmg"}：CCEMG；\texttt{"dmg"}：去均值 MG & 14\\
\addlinespace[1pt]
\texttt{fe <- plm(y \textasciitilde{} x1 + x2, data = pd, model = "within")}\newline \texttt{pcdtest(fe, test = "cd")} & Pesaran CD 检验 & \texttt{test = "lm"}（BP LM）、\texttt{"sclm"}（$CD^*$，标准化 LM）、\texttt{"bcsclm"}（偏误校正的标准化 LM） & 14\\
\end{longtable}}


\begin{note}[\texttt{plm()} 没有 PCA 选项]
\texttt{plm(..., model = "PCA")} 是不存在的写法。\texttt{plm()} 的 \texttt{model} 参数可取 \texttt{"within"}、\texttt{"random"}、\texttt{"pooling"}、\texttt{"between"}、\texttt{"fd"}（以及 Hausman–Taylor 用的 \texttt{"ht"}），都不估计因子结构。因子模型要用上表的 \texttt{eigen} 手工计算或用 phtt 等专门的包。
\end{note}

\begin{english}
\begin{itemize}
  \item We test for cross-sectional dependence in the FE residuals with the Pesaran CD test, which strongly rejects the null of cross-sectional independence.
  \item To account for unobserved common factors, we estimate the model by the CCE mean group estimator and by Bai's (2009) interactive fixed effects estimator.
  \item Driscoll--Kraay standard errors are robust to cross-sectional and serial correlation, but they do not correct the bias that arises when the factors are correlated with the regressors.
\end{itemize}
\end{english}

%======================================================================
\section{DiD、事件研究与交错处理（第 15 章）}\label{sec:UappB-did}

\subsection{\texorpdfstring{$2\times 2$}{2x2} DiD、TWFE 与异质效应}

{\renewcommand{\arraystretch}{1.12}\begin{longtable}{@{}>{\raggedright\arraybackslash\small}p{5.3cm}>{\raggedright\arraybackslash\small}p{3.8cm}>{\raggedright\arraybackslash\small}p{5.3cm}>{\centering\arraybackslash\small}p{0.9cm}@{}}
\toprule
\textbf{命令} & \textbf{作用} & \textbf{常用选项 / 注意} & \textbf{章}\\
\midrule
\endfirsthead
\toprule
\textbf{命令} & \textbf{作用} & \textbf{常用选项 / 注意} & \textbf{章}\\
\midrule
\endhead
\bottomrule
\endlastfoot

\texttt{reg y i.d\#\#i.w, vce(cluster id)} & $2\times2$ DiD 回归（$d$、$w$ 与交互） & 交互项 \texttt{1.d\#1.w} 的系数即 DiD；\texttt{d.f} 是差分算子，\key{不是}交互 & 15\\
\addlinespace[1pt]
\texttt{xtset id year}\newline \texttt{xtreg y D i.year, fe vce(cluster id)} & TWFE（$D=d\times w$） & 共同时点下 = 四格 DiD & 15\\
\addlinespace[1pt]
\texttt{reghdfe y D, absorb(id year) vce(cluster id)}\newline {\footnotesize\rmfamily [SSC]} & 同上（高维 FE） & — & 15\\
\addlinespace[1pt]
\texttt{xtreg y D x i.year, fe vce(cluster id)} & 加时变协变量的 TWFE & 慎用：协变量可能受处理影响（坏控制）；见 Caetano、Callaway 等的讨论 & 15\\
\addlinespace[1pt]
\texttt{gen D3 = d*(year==3)}\newline \texttt{gen D4 = d*(year==4)}\newline \texttt{gen D5 = d*(year==5)}\newline \texttt{xtreg y D3 D4 D5 i.year, fe vce(cluster id)} & 时变处理效应 $\tau_t$ & 先生成各处理期的处理虚拟变量再回归 & 15\\
\addlinespace[1pt]
\texttt{reghdfe y i.id\#c.D, absorb(id year) vce(robust)} & 个体异质效应 $\beta_i$（每个处理单位一个交互项） & 2025 年真题第 2(b) 题（见第 17 章）；控制单位 $D=0$，其交互项自动略去 & 15\\
\addlinespace[1pt]
\texttt{test \_b[1.id\#c.D] = \_b[2.id\#c.D] = \_b[3.id\#c.D]} & 检验处理效应同质 $H_0:\beta_1=\beta_2=\beta_3$ & — & 15\\
\end{longtable}}


\subsection{事件研究与交错处理（异质稳健估计量）}

\subsubsection{事件研究回归、事前趋势检验与事件研究图}

{\renewcommand{\arraystretch}{1.12}\begin{longtable}{@{}>{\raggedright\arraybackslash\small}p{5.3cm}>{\raggedright\arraybackslash\small}p{3.8cm}>{\raggedright\arraybackslash\small}p{5.3cm}>{\centering\arraybackslash\small}p{0.9cm}@{}}
\toprule
\textbf{命令} & \textbf{作用} & \textbf{常用选项 / 注意} & \textbf{章}\\
\midrule
\endfirsthead
\toprule
\textbf{命令} & \textbf{作用} & \textbf{常用选项 / 注意} & \textbf{章}\\
\midrule
\endhead
\bottomrule
\endlastfoot

\texttt{gen rel = year - first\_treat}\newline 生成 \texttt{lead2}–\texttt{lead4}、\texttt{lag0}–\texttt{lag3}（见模板三） & 事件时间与前导/滞后虚拟变量 & 略去 $e=-1$（参照期） & 15\\
\addlinespace[1pt]
\texttt{reghdfe y lead4 lead3 lead2 lag0 lag1 lag2 lag3, absorb(id year) vce(cluster id)} & 事件研究回归 & 共同时点下各系数为相对 $e=-1$ 的 DiD & 15\\
\addlinespace[1pt]
\texttt{test lead4 lead3 lead2} & 事前趋势联合检验 & 不拒绝 $\neq$ 平行趋势成立 & 15\\
\addlinespace[1pt]
\texttt{coefplot, keep(lead* lag*) vertical yline(0)}\newline {\footnotesize\rmfamily [SSC]} & 事件研究图 & — & 15\\
\end{longtable}}


\subsubsection{扩展 TWFE 与普通 TWFE 的诊断}

{\renewcommand{\arraystretch}{1.12}\begin{longtable}{@{}>{\raggedright\arraybackslash\small}p{5.3cm}>{\raggedright\arraybackslash\small}p{3.8cm}>{\raggedright\arraybackslash\small}p{5.3cm}>{\centering\arraybackslash\small}p{0.9cm}@{}}
\toprule
\textbf{命令} & \textbf{作用} & \textbf{常用选项 / 注意} & \textbf{章}\\
\midrule
\endfirsthead
\toprule
\textbf{命令} & \textbf{作用} & \textbf{常用选项 / 注意} & \textbf{章}\\
\midrule
\endhead
\bottomrule
\endlastfoot

\texttt{reghdfe y tau\_*, absorb(id year) vce(cluster id)} & 手工扩展 TWFE（每个已处理的 cohort$\times$期一个虚拟变量） & \texttt{tau\_r\_s = (first\_treat==r)*(year==s)}，$s\ge r$ & 15\\
\addlinespace[1pt]
\texttt{jwdid y, ivar(id) tvar(year) gvar(first\_treat)}\newline {\footnotesize\rmfamily [SSC]} & Wooldridge 扩展 TWFE（Rios-Avila） & 之后 \texttt{estat simple}（总 ATT）、\texttt{estat event}、\texttt{estat group}；另一个实现扩展 TWFE 的用户命令是 \texttt{wooldid} & 15\\
\addlinespace[1pt]
\texttt{xthdidregress twfe (y) (D), group(id)}\newline \texttt{xthdidregress aipw (y x) (D x), group(id) controlgroup(notyet)}\newline {\footnotesize\rmfamily [官方，Stata 18 起]} & 官方异质 DiD：\texttt{twfe}、\texttt{ra}、\texttt{ipw}、\texttt{aipw} 四种方法 & 第一个括号为结果模型，第二个为处理模型；\texttt{group()} 为处理发生的层级；之后 \texttt{estat aggregation, dynamic graph} & 15\\
\addlinespace[1pt]
\texttt{bacondecomp y D, ddetail}\newline {\footnotesize\rmfamily [SSC]} & Goodman-Bacon 分解 & 诊断普通 TWFE & 15\\
\addlinespace[1pt]
\texttt{twowayfeweights y id year D, type(feTR)}\newline {\footnotesize\rmfamily [SSC]} & dCDH 权重（负权重诊断） & — & 15\\
\end{longtable}}


\subsubsection{异质稳健估计量：CS、SA、BJS 与 dCDH}

{\renewcommand{\arraystretch}{1.12}\begin{longtable}{@{}>{\raggedright\arraybackslash\small}p{5.3cm}>{\raggedright\arraybackslash\small}p{3.8cm}>{\raggedright\arraybackslash\small}p{5.3cm}>{\centering\arraybackslash\small}p{0.9cm}@{}}
\toprule
\textbf{命令} & \textbf{作用} & \textbf{常用选项 / 注意} & \textbf{章}\\
\midrule
\endfirsthead
\toprule
\textbf{命令} & \textbf{作用} & \textbf{常用选项 / 注意} & \textbf{章}\\
\midrule
\endhead
\bottomrule
\endlastfoot

\texttt{csdid y, ivar(id) time(year) gvar(first\_treat) notyet}\newline {\footnotesize\rmfamily [SSC]} & Callaway–Sant'Anna $ATT(g,t)$ & \texttt{notyet}：用尚未处理组作控制；之后 \texttt{estat event}、\texttt{csdid\_plot} & 15\\
\addlinespace[1pt]
\texttt{eventstudyinteract y lead* lag*, cohort(first\_treat) control\_cohort(never) absorb(i.id i.year) vce(cluster id)}\newline {\footnotesize\rmfamily [SSC]} & Sun–Abraham 交互加权事件研究 & \texttt{never}：从未处理的指示变量 & 15\\
\addlinespace[1pt]
\texttt{did\_imputation y id year first\_treat, horizons(0/5) pretrends(3)}\newline {\footnotesize\rmfamily [SSC]} & Borusyak–Jaravel–Spiess 插补估计量 & 事前检验单独用未处理观测 & 15\\
\addlinespace[1pt]
\texttt{did\_multiplegt\_dyn y id year D, effects(5) placebo(3)}\newline {\footnotesize\rmfamily [SSC]} & de Chaisemartin–D'Haultf\oe uille 动态效应 & 允许处理开关；旧版命令为 \texttt{did\_multiplegt} & 15\\
\end{longtable}}


\begin{english}
\begin{itemize}
  \item The DiD estimate is the coefficient on $D_{it}=d_iw_t$ in a regression with unit and time fixed effects, and the standard errors are clustered by unit.
  \item In the event-study specification the period $e=-1$ is omitted as the reference period; jointly insignificant leads support, but do not prove, the parallel trends assumption.
  \item With staggered adoption and heterogeneous effects, we complement the static TWFE estimate with heterogeneity-robust estimators (Callaway--Sant'Anna, Sun--Abraham, and the imputation estimator).
\end{itemize}
\end{english}

%======================================================================
\section{合成控制、面板数据法与矩阵补全（第 16 章）}\label{sec:UappB-scm}

\subsection{Stata 与 R 命令：SCM、安慰剂推断、回归控制法与矩阵补全}

{\renewcommand{\arraystretch}{1.12}\begin{longtable}{@{}>{\raggedright\arraybackslash\small}p{5.3cm}>{\raggedright\arraybackslash\small}p{3.8cm}>{\raggedright\arraybackslash\small}p{5.3cm}>{\centering\arraybackslash\small}p{0.9cm}@{}}
\toprule
\textbf{命令} & \textbf{作用} & \textbf{常用选项 / 注意} & \textbf{章}\\
\midrule
\endfirsthead
\toprule
\textbf{命令} & \textbf{作用} & \textbf{常用选项 / 注意} & \textbf{章}\\
\midrule
\endhead
\bottomrule
\endlastfoot

\texttt{tsset id year}\newline \texttt{synth y x1 x2 y(1985) y(1988) y(1990), trunit(1) trperiod(1991) fig}\newline {\footnotesize\rmfamily [SSC]} & 合成控制法（ADH） & \texttt{trunit()} 处理单位，\texttt{trperiod()} 处理开始期；\texttt{y(1985)} 等为作预测变量的处理前结果 & 16\\
\addlinespace[1pt]
\texttt{synth\_runner y x1 x2 y(1985) y(1988) y(1990), trunit(1) trperiod(1991) gen\_vars}\newline {\footnotesize\rmfamily [SSC]} & 空间安慰剂推断（秩 $p$ 值） & — & 16\\
\addlinespace[1pt]
\texttt{xtset id year}\newline \texttt{rcm y, trunit(1) trperiod(2004) method(lasso)}\newline {\footnotesize\rmfamily [SSC]} & 回归控制法 / 面板数据法（PDA） & Yan 与 Chen 编写；\texttt{method()} 指定挑选控制单位的方法（如 \texttt{lasso}）；完整语法见 \texttt{help rcm} & 16\\
\addlinespace[1pt]
R：\texttt{Synth}、\texttt{tidysynth}、\texttt{augsynth} & SCM；增广 SCM & — & 16\\
\addlinespace[1pt]
R：\texttt{pampe} & PDA & — & 16\\
\addlinespace[1pt]
R：\texttt{MCPanel}、\texttt{gsynth}、\texttt{fect} & 矩阵补全、广义合成控制 & — & 16\\
\end{longtable}}


\begin{english}
\begin{itemize}
  \item The synthetic control is a convex combination of control units whose weights are chosen to match the treated unit's pre-treatment outcomes and predictors.
  \item Inference is based on in-space placebo tests: the $p$-value is the rank of the treated unit's post/pre RMSPE ratio among all units divided by the number of units.
  \item The panel data approach (regression control method) instead uses an unrestricted linear combination of control units, estimated by OLS on the pre-treatment periods.
\end{itemize}
\end{english}

%======================================================================
\section{R 命令对照：fixest、plm、did 等（拓展）}\label{sec:UappB-r}

\subsection{与 Stata 命令对应的 R 写法}

正文主要用 Stata；下表给出常用 R 包中对应的标准写法。\texttt{df} 为数据框，\texttt{pd} 为 \texttt{pdata.frame} 对象。

{\renewcommand{\arraystretch}{1.12}\begin{longtable}{@{}>{\raggedright\arraybackslash\small}p{5.3cm}>{\raggedright\arraybackslash\small}p{3.8cm}>{\raggedright\arraybackslash\small}p{5.3cm}>{\centering\arraybackslash\small}p{0.9cm}@{}}
\toprule
\textbf{命令} & \textbf{作用} & \textbf{常用选项 / 注意} & \textbf{章}\\
\midrule
\endfirsthead
\toprule
\textbf{命令} & \textbf{作用} & \textbf{常用选项 / 注意} & \textbf{章}\\
\midrule
\endhead
\bottomrule
\endlastfoot

\texttt{m <- lm(y \textasciitilde{} x1 + x2, data = df)}\newline \texttt{lmtest::coeftest(m, vcov = sandwich::vcovHC(m, type = "HC1"))} & OLS + HC1 稳健 se & \texttt{sandwich::vcovCL(m, cluster = \textasciitilde{}id)} 为聚类 se & 5\\
\addlinespace[1pt]
\texttt{AER::ivreg(y \textasciitilde{} x1 + x2 | z1 + z2 + x2, data = df)} & 2SLS（竖线后为全部工具） & 外生回归元也要写进工具 & 8\\
\addlinespace[1pt]
\texttt{fixest::feols(y \textasciitilde{} x | id + year, data = df, cluster = \textasciitilde{}id)} & TWFE（个体 + 年份 FE），聚类 se & 竖线后为固定效应 & 10, 15\\
\addlinespace[1pt]
\texttt{feols(y \textasciitilde{} x2 | id | x1 \textasciitilde{} z1 + z2, data = df)} & 含 FE 的 2SLS & 第三段为“内生变量 ~ 工具” & 8, 10\\
\addlinespace[1pt]
\texttt{feols(y \textasciitilde{} sunab(first\_treat, year) | id + year, data = df)} & Sun–Abraham 事件研究 & \texttt{fixest::sunab()} & 15\\
\addlinespace[1pt]
\texttt{plm(y \textasciitilde{} x, data = pd, model = "within")} & FE（WG） & 其他取值：\texttt{"random"}、\texttt{"pooling"}、\texttt{"between"}、\texttt{"fd"} & 10–11\\
\addlinespace[1pt]
\texttt{phtest(fe, re)} & Hausman 检验（plm） & \texttt{fe}、\texttt{re} 为两个 plm 估计结果 & 11\\
\addlinespace[1pt]
\texttt{lmtest::coeftest(fe, vcov = vcovHC(fe, cluster = "group"))} & FE 的按个体聚类 se & 对 plm 对象默认用 Arellano 方法 & 10\\
\addlinespace[1pt]
\texttt{pgmm(y \textasciitilde{} lag(y, 1) | lag(y, 2:99), data = pd, effect = "individual", model = "twosteps")} & AB 差分 GMM（plm） & \texttt{transformation = "ld"} 为系统 GMM；参数见 \texttt{?pgmm} & 12\\
\addlinespace[1pt]
\texttt{out <- did::att\_gt(yname = "y", tname = "year", idname = "id", gname = "first\_treat", data = df)}\newline \texttt{did::aggte(out, type = "dynamic")} & Callaway–Sant'Anna $ATT(g,t)$ 及事件时间聚合 & R 包 did；从未处理组的 \texttt{first\_treat} 记 0 & 15\\
\addlinespace[1pt]
R：\texttt{didimputation}、\texttt{HonestDiD} & BJS 插补；平行趋势敏感性分析 & — & 15\\
\end{longtable}}


\begin{english}
\begin{itemize}
  \item In R, two-way fixed effects models with clustered standard errors can be estimated with \texttt{fixest::feols}, and panel GMM estimators with \texttt{plm::pgmm}.
  \item The Callaway--Sant'Anna group-time average treatment effects are estimated with \texttt{did::att\_gt} and aggregated by event time with \texttt{did::aggte}.
\end{itemize}
\end{english}

%======================================================================
\section{完整实证流程模板与常见陷阱}\label{sec:UappB-templates}

\subsection{模板一：IV 的三大检验流程（第 9 章）}

\begin{lstlisting}
* (0) 基准 OLS
regress y x1 x2, vce(robust)
* (1) 2SLS: x1 内生, z1 z2 为排除工具, x2 外生; first = 打印第一阶段
ivregress 2sls y (x1 = z1 z2) x2, vce(robust) first
estat firststage      // 检验一: 相关性 (F, 偏 R2)
estat overid          // 检验二: 过度识别 (L-K = 排除工具 2 - 内生变量 1 = 1)
estat endogenous      // 检验三: 内生性
* (2) 两步有效 GMM; 加 igmm 为迭代 GMM
ivregress gmm y (x1 = z1 z2) x2, wmatrix(robust)
estat overid          // Hansen J
* (3) LIML (弱工具时的替代)
ivregress liml y (x1 = z1 z2) x2, vce(robust)
\end{lstlisting}

\subsection{模板二：动态面板 GMM 流程（第 12 章）}

\begin{lstlisting}
ssc install xtabond2
xtset id year
tab year, gen(yr)                                    // 年份虚拟变量
reg   y L.y x yr*, vce(cluster id)                   // 混合 OLS: rho 上偏
xtreg y L.y x yr*, fe vce(cluster id)                // WG: rho 下偏 (Nickell)
xtabond2 y L.y x yr*, gmm(L.y) iv(x yr*) noleveleq twostep robust             // AB
xtabond2 y L.y x yr*, gmm(L.y) iv(x yr*) twostep robust                       // System
xtabond2 y L.y x yr*, gmm(L.y) iv(x yr*) noleveleq orthogonal twostep robust  // FOD
xtabond2 y L.y x yr*, gmm(L.y, laglimits(1 3) collapse) iv(x yr*) twostep robust // 限制工具
* 看输出: AR(1) 应显著, AR(2) 不显著; Hansen J 不拒绝; 工具数 < 个体数
\end{lstlisting}

\subsection{模板三：DiD、事件研究与交错处理（第 15 章）}

\begin{lstlisting}
xtset id year
xtreg y D i.year, fe vce(cluster id)              // 共同时点 TWFE = DiD
* first_treat: 首次处理年份; 从未处理者 csdid/jwdid 要求记 0, did_imputation 要求记缺失
gen rel = year - first_treat                      // 事件时间 (从未处理者的 lead/lag 均为 0)
forvalues k = 2/4 {
    gen lead`k' = (rel == -`k')                   // e = -1 略去
}
forvalues k = 0/3 {
    gen lag`k' = (rel == `k')
}
reghdfe y lead4 lead3 lead2 lag0 lag1 lag2 lag3, absorb(id year) vce(cluster id)
test lead4 lead3 lead2                            // 事前趋势检验
bacondecomp y D, ddetail                          // 交错处理时诊断普通 TWFE
jwdid y, ivar(id) tvar(year) gvar(first_treat)    // 扩展 TWFE
estat event
csdid y, ivar(id) time(year) gvar(first_treat) notyet
estat event
\end{lstlisting}

\subsection{常见语法与使用陷阱清单}

\begin{enumerate}
  \item \key{交互项}：\texttt{d.f} 是差分算子 \texttt{D.} 作用于 \texttt{f}，不是 $d\times f$；交互写 \texttt{i.d\#i.f} 或 \texttt{c.d\#c.f}，含主效应写 \texttt{i.d\#\#i.f}（第 15 章）。
  \item \key{FE 的标准误}：\texttt{xtreg, fe vce(robust)} 实际给出按个体聚类的 se；手工去均值后 \texttt{reg ydot xdot} 的经典 se 自由度错（用 $NT-K$ 而非 $NT-N-K$），需乘 $\sqrt{(NT-K)/(NT-N-K)}$（第 10 章）。
  \item \key{Hausman 检验}：\texttt{hausman} 的顺序是“先一致的、后有效的”；加 \texttt{sigmamore} 避免方差之差非正定；它只在同方差下有效，稳健版本用 Mundlak 回归 + \texttt{test xbar}（第 11 章）。
  \item \key{2SLS}：不要手工两步做（se 错）；\texttt{estat firststage} 的 Stock–Yogo 临界值只在非稳健设定下报告（第 9 章）。
  \item \key{聚类}：簇数 $G$ 小（$<30$）时 CRVE 不可靠；Stata 自动用 $t_{G-1}$；$G=1$ 时 CRVE 为 0（第 5 章）。
  \item \key{\texttt{xtabond2} 的三种估计量}：\texttt{noleveleq} $\to$ AB，默认 $\to$ System，\texttt{noleveleq orthogonal} $\to$ FOD；工具过多会使 Hansen $J$ 的 $p$ 值虚高，用 \texttt{collapse}、\texttt{laglimits()} 控制并报告工具数（第 12 章）。
  \item \key{PCA 软件约定}：统计软件把变量放在列上；MATLAB/Python 的 \texttt{eig}/\texttt{eigh} 按特征值升序返回，R 的 \texttt{eigen} 按降序返回（第 13 章）；\texttt{plm()} 没有 \texttt{model = "PCA"}（第 14 章）。
  \item \key{事件研究}：必须略去一个事件时间（通常 $e=-1$），否则事件时间虚拟变量之和与个体固定效应完全共线（第 15 章）。
  \item \key{多个约束}：连写两条 \texttt{test} 是两个单独检验；联合检验要把约束放进同一条 \texttt{test} 的括号里，或在第二条加 \texttt{accumulate}（第 6 章）。
\end{enumerate}

\begin{english}
\begin{itemize}
  \item First, we estimate the baseline model by OLS with robust standard errors; second, we estimate it by 2SLS and report the first-stage $F$ statistic, the overidentification test and the endogeneity test; finally, we check robustness with efficient GMM and LIML.
  \item For the dynamic panel, we report pooled OLS and within estimates as upper and lower bounds for $\rho$, and verify that the GMM estimate lies between them.
  \item For the policy evaluation, we report the TWFE estimate, an event-study plot with $e=-1$ as the reference period, and heterogeneity-robust estimates under staggered adoption.
\end{itemize}
\end{english}
